% Generated draft: requires mathematical and source review.
\documentclass[11pt]{article}
\usepackage[margin=25mm]{geometry}
\usepackage{fontspec}
\setmainfont{Latin Modern Roman}
\usepackage{amsmath,amssymb,mathtools,unicode-math}
\setmathfont{Latin Modern Math}
\usepackage{enumitem,needspace,xurl,hyperref}
\setlength{\parindent}{0pt}
\setlength{\parskip}{5pt}
\setlength{\emergencystretch}{3em}
\newcommand{\N}{\mathbb{N}}
\newcommand{\Z}{\mathbb{Z}}
\newcommand{\Q}{\mathbb{Q}}
\newcommand{\R}{\mathbb{R}}
\newcommand{\C}{\mathbb{C}}
\newcommand{\F}{\mathbb{F}}
\newcommand{\A}{\mathbb{A}}
\newcommand{\PP}{\mathbb{P}}
\newcommand{\E}{\mathbb{E}}
\newcommand{\eps}{\varepsilon}
\newcommand{\dd}{\,\mathrm{d}}
\begin{document}
\section*{UnsolvedMath ID: 20000431}
\textbf{Source problem code:} catalogue entry
\subsection{Definitions and mathematical statement}
\textbf{Original question (M. Flach).} Find a natural definition of Weil--étale topology for which integral cohomology vanishes in degrees greater than three. Flach comments that the then-current definition has the unwanted pattern of vanishing in odd degrees and infinitely generated groups in even degrees, and that it is unclear whether a better definition is near or truncation of the cohomology complex is the only solution. The low-degree range and underlying arithmetic object are not specified in the extracted sentence; the attached aid specializes to $\operatorname{Spec}\Z$. The historical odd-degree comment is not asserted as a newly checked statement about every degree or every modern theory.

\textbf{Attached unreviewed necessary test.} Let a proposed compactified Grothendieck topos for $\operatorname{Spec}\Z$ have ordinary and compactly supported derived sections $R\Gamma$ and $R\Gamma_c$. The aid imposes the following explicit hypotheses: (H1) retain the indicated Lichtenbaum integral groups in degrees zero through three, namely ordinary $H^0=\Z$, $H^1=H^2=0$, $H^3=\Z/2$, and compactly supported $H_c^0=H_c^1=H_c^2=0$, $H_c^3=\Z/2$; (H2) both theories vanish in negative degrees and degrees greater than three; (H3) both derived section theories respect the constant-sheaf coefficient sequence
\[0\longrightarrow\Z\xrightarrow{m}\Z\longrightarrow\Z/m\longrightarrow0\qquad(m\ge1).\]
The compactification, support functor and claimed historical verification of these low groups require independent review.

Under (H1)--(H3), the reported necessary derived forms are, in $D(\Z)$,
\[R\Gamma(\Z)\simeq\Z\oplus(\Z/2)[-3],\qquad R\Gamma_c(\Z)\simeq(\Z/2)[-3].\]
These are derived isomorphism statements, not an asserted canonical splitting or construction of a topology. Here a group shifted by $[-3]$ lies in cohomological degree three.

The claimed finite-coefficient fingerprint is: for every $m\ge1$, if $m$ is even then
\[H_c^i(\Z/m)\cong\begin{cases}\Z/2&i=2,3,\\0&\text{otherwise},\end{cases}\]
and if $m$ is odd all $H_c^i(\Z/m)$ vanish. Ordinary cohomology has the same degree-two and degree-three parity pattern and additionally $H^0(\Z/m)=\Z/m$, with all other degrees zero. The degree-two term is the coefficient-sequence/Bockstein effect and cannot be omitted merely because integral $H^2$ vanishes.

The aid describes the requested bounded integral topology as still open. Its parity test is only a necessary falsification criterion for theories satisfying the stated low-degree and coefficient hypotheses; it neither supplies the missing topology nor proves those hypotheses are the unique desirable axioms.
\subsection{Short English statement}
Flach asks for a natural Weil--étale topology with integral cohomology bounded above by degree three. The aid supplies only a necessary test on $\operatorname{Spec}\Z$: preserving specified low groups and coefficient exactness forces finite-coefficient groups in degrees two and three precisely for even moduli. It does not construct the topology.
\subsection{Sources}
\begin{itemize}
\item Local UnsolvedMath registry, permanent ID 20000431, source code catalogue entry; \texttt{aim-algebraic-number-theory-notes.json}, notes[82], tag \texttt{problem}.
\item AIM workshop \emph{The Tate conjecture}, M. Flach's question: \url{https://aimath.org/WWN/tateconjecture/tateconjecture.pdf}. Extraction/source-item metadata say 0; supplied verification identifies original TeX Problem 20 and describes the OCR damage. No PDF or TeX source was consulted here.
\item Attached research attempt 1 attributes the low-degree groups to Lichtenbaum without a publication locator. Result \texttt{partial\_result}, problem status at run \texttt{open}, rigor label \texttt{proved}, novelty confidence low. It is a machine-generated, non-peer-reviewed aid requiring independent verification before citation. Original literature field empty.
\end{itemize}
\end{document}
